% !TeX encoding = UTF-8
\documentclass[variant=guia,cover=institutional,mode=print]{ua-engineering-v6-article}
% !TeX encoding = UTF-8
\UASetUniversity{Universidad Austral}
\UASetFaculty{Facultad de Ingeniería}
\UASetDepartment{Departamento de Computación}
\UASetCampus{Pilar}
\UASetCareer{Ingeniería Informática}
\UASetCommission{Comisión A}
\UASetProfessor{Equipo docente}
\UASetSemester{Segundo cuatrimestre 2026}
\UASetCourse{Sistemas Distribuidos}
\UASetDate{2026}

\UASetSemester{Segundo cuatrimestre 2026}
\UASetDate{2026}
\UASetClassNumber{09}
\UASetClassTitle{Arquitecturas distribuidas reales}
\title{Clase 09 --- Ficha de aula: fronteras, ownership, outbox y sagas}
\author{\UAProfessor}
\date{\UADate}
\begin{document}
\maketitle
\section{1. Intuiciones iniciales}
\begin{tabular}{p{0.20\textwidth}p{0.32\textwidth}p{0.34\textwidth}}
\toprule
Concepto & Imagen mental & Definición técnica al final \\
\midrule
Monolito modular & edificio/ habitaciones & \rule{0pt}{1cm} \\
Microservicio & edificio/ caminos & \rule{0pt}{1cm} \\
Bounded context & idioma local & \rule{0pt}{1cm} \\
Ownership & escritura del dato & \rule{0pt}{1cm} \\
RPC & llamada telefónica & \rule{0pt}{1cm} \\
Evento & aviso durable & \rule{0pt}{1cm} \\
Outbox & copia en mismo libro & \rule{0pt}{1cm} \\
Saga & itinerario compensable & \rule{0pt}{1cm} \\
\bottomrule
\end{tabular}

\newpage

\begin{uainfo}
\textbf{Precisión para la ficha:} un bounded context no obliga a crear un proceso. Primero se justifica la frontera semántica; después se decide si la autonomía ganada paga una frontera de red.
\end{uainfo}

\section{2. Canvas de frontera para Asteria}
Elija una capacidad: combate, inventario, mercado, pagos, ranking o chat.

\begin{tabular}{p{0.30\textwidth}p{0.56\textwidth}}
\toprule
Pregunta & Respuesta \\
\midrule
Capacidad y lenguaje & \rule{0pt}{1cm} \\
Invariante local & \rule{0pt}{1cm} \\
Owner/equipo & \rule{0pt}{1cm} \\
Dato exclusivo & \rule{0pt}{1cm} \\
Presión para separar & \rule{0pt}{1cm} \\
Falla nueva & \rule{0pt}{1cm} \\
SLO / métrica & \rule{0pt}{1cm} \\
Alternativa simple & \rule{0pt}{1cm} \\
\bottomrule
\end{tabular}

\newpage
\section{3. Prima distribuida}
Transforme una llamada local en RPC. Complete costos:

\begin{center}
\begin{tikzpicture}[node distance=3cm,>=Latex,every node/.style={align=center}]
\node[draw=UAIngOrange,rounded corners,thick,text width=3cm] (a) {Módulo A};
\node[draw=UAIngNavy,rounded corners,thick,right=of a,text width=3cm] (b) {Servicio B};
\draw[->,very thick] (a)--node[above]{red}(b);
\end{tikzpicture}
\end{center}


\begin{itemize}
\item latencia: \hrulefill
\item timeout: \hrulefill
\item retry: \hrulefill
\item contrato: \hrulefill
\item seguridad: \hrulefill
\item tracing: \hrulefill
\item despliegue: \hrulefill
\item on-call: \hrulefill
\end{itemize}


Calcule $0{,}999^4$: \hrulefill. ¿Por qué es optimista? \vspace{1cm}

\newpage
\section{4. RPC, evento o consulta}

\begin{tabular}{p{0.28\textwidth}p{0.18\textwidth}p{0.44\textwidth}}
\toprule
Interacción & Elección & Justificación y falla \\
\midrule
Validar pago & & \rule{0pt}{1.1cm} \\
Enviar email & & \rule{0pt}{1.1cm} \\
Actualizar ranking & & \rule{0pt}{1.1cm} \\
Consultar perfil & & \rule{0pt}{1.1cm} \\
Reservar objeto & & \rule{0pt}{1.1cm} \\
\bottomrule
\end{tabular}

Clasifique:
\begin{itemize}
\item \texttt{ReserveSeat}: \hrulefill
\item \texttt{SeatReserved}: \hrulefill
\item \texttt{GetSeatAvailability}: \hrulefill
\end{itemize}

\newpage
\section{5. Dual write y outbox}
Dibuje los dos órdenes y marque la ventana de crash.
\vspace{4cm}

Diseñe outbox:
\begin{itemize}
\item event ID: \hrulefill
\item aggregate ID y sequence: \hrulefill
\item estado del relay: \hrulefill
\item deduplicación: \hrulefill
\item métricas: \hrulefill
\end{itemize}

\newpage
\section{6. Saga de compra}
Complete estados, retries y compensaciones.

\begin{center}
\begin{tikzpicture}[node distance=1cm,>=Latex,every node/.style={align=center}]
\node[draw=UAIngOrange,rounded corners,thick] (a) {Reservar};
\node[draw=UAIngNavy,rounded corners,thick,right=of a] (b) {Pagar};
\node[draw=UAIngOrange,rounded corners,thick,right=of b] (c) {Transferir};
\node[draw=UAIngNavy,rounded corners,thick,right=of c] (d) {Notificar};
\draw[->,thick] (a)--(b)--(c)--(d);
\end{tikzpicture}
\end{center}

\begin{tabular}{p{0.20\textwidth}p{0.20\textwidth}p{0.24\textwidth}p{0.18\textwidth}}
\toprule
Paso & Falla & Acción & ¿Rollback real? \\
\midrule
Reserva & \rule{0pt}{1cm} & & \\
Pago & \rule{0pt}{1cm} & & \\
Transferencia & \rule{0pt}{1cm} & & \\
Email & \rule{0pt}{1cm} & & \\
\bottomrule
\end{tabular}

\newpage
\section{7. Migración, métricas y exit ticket}
Proponga una ruta Strangler:
\begin{enumerate}
\item capacidad inicial: \hrulefill
\item router/proxy: \hrulefill
\item fuente de verdad durante coexistencia: \hrulefill
\item rollback: \hrulefill
\item métrica que justifica continuar: \hrulefill
\end{enumerate}

\subsection*{Exit ticket}
\begin{enumerate}
\item Diferencia entre monolito modular y monolito distribuido.
\item Presión concreta que justifica una extracción.
\item Qué resuelve outbox y qué deja abierto.
\item Por qué una compensación no es rollback.
\end{enumerate}
\end{document}
